%% ISMB / Bioinformatics (OUP) submission template
%% Title: Systematic benchmarking of functional module partitioning in GO pathway subnetworks
%% Compile with: pdflatex -> bibtex -> pdflatex -> pdflatex
\documentclass[twocolumn,10pt,letterpaper]{article}

\usepackage[T1]{fontenc}
\usepackage{newtxtext}
\usepackage{graphicx}
\usepackage{amsmath,amssymb}
\usepackage{booktabs}
\usepackage{multirow}
\usepackage[numbers,sort&compress]{natbib}
\usepackage[colorlinks=true,linkcolor=blue,citecolor=blue,urlcolor=blue]{hyperref}
\usepackage{float}
\usepackage[letterpaper,top=1in,bottom=1in,left=1in,right=1in,columnsep=0.25in,footskip=0.35in]{geometry}
\setlength{\floatsep}{5pt plus 1pt minus 1pt}
\setlength{\textfloatsep}{7pt plus 1pt minus 1pt}
\setlength{\intextsep}{5pt plus 1pt minus 1pt}
\setlength{\abovedisplayskip}{4pt plus 1pt minus 1pt}
\setlength{\belowdisplayskip}{4pt plus 1pt minus 1pt}
\setlength{\abovedisplayshortskip}{2pt}
\setlength{\belowdisplayshortskip}{2pt}
\usepackage{caption}
\captionsetup[figure]{font=footnotesize,labelfont=bf,skip=3.5pt,name={Fig.},labelsep=period}
\captionsetup[table]{font=footnotesize,labelfont=bf,skip=3.5pt,name={Table},labelsep=period}

% ---- 排版微调：避免溢出、平衡字间距 ----
\setlength{\emergencystretch}{2.5em}
\tolerance=1200
\widowpenalty=1000
\clubpenalty=1000
\setlength{\parskip}{0pt}
\setlength{\abovedisplayskip}{4pt}
\setlength{\belowdisplayskip}{4pt}
\usepackage{titlesec}
\titleformat{\section}[hang]{\normalfont\large\bfseries}{\thesection}{0.6em}{}
\titleformat{\subsection}[hang]{\normalfont\normalsize\bfseries}{\thesubsection}{0.6em}{}
\titleformat{\subsubsection}[runin]{\normalfont\small\bfseries}{}{0em}{}[.]
\titlespacing*{\section}{0pt}{6pt plus 2pt minus 1pt}{2.5pt}
\titlespacing*{\subsection}{0pt}{4pt plus 1pt minus 1pt}{1.5pt}
\titlespacing*{\subsubsection}{0pt}{4pt}{0.5em}

%% ---- 跨栏标题 + 摘要（ISMB/OUP 样式）----
\renewcommand{\topfraction}{0.95}
\renewcommand{\bottomfraction}{0.9}
\renewcommand{\textfraction}{0.06}
\renewcommand{\floatpagefraction}{0.72}
\begin{document}

\twocolumn[{
\begin{center}
{\LARGE\bfseries FACEO-R: Functionally Anchored Overlapping Community Detection for
Robust Pathway Module Discovery\par}
\vspace{0.85em}
{\large Mingyu Sun\textsuperscript{2} \quad Weixiong Zhang\textsuperscript{1,\,2,\,3,\,*}\par}
\vspace{0.28em}
{\footnotesize\textsuperscript{1}Hong Kong Jockey Club STEM Laboratory for Genomics and AI
in Healthcare, Hong Kong, China\par}
{\footnotesize\textsuperscript{2}Department of Health Technology and Informatics, The Hong
Kong Polytechnic University, Hong Kong, China\par}
{\footnotesize\textsuperscript{3}Department of Data Science and Artificial Intelligence,
The Hong Kong Polytechnic University, Hong Kong, China\par}
{\footnotesize\textsuperscript{*}Correspondence: weixiong.zhang@polyu.edu.hk;
mingyu.sun@connect.polyu.hk\par}
\end{center}
\vspace{0.35em}
\hrule height 0.4pt
\vspace{0.60em}
\begin{minipage}{0.955\textwidth}
\footnotesize
\setlength{\parindent}{0pt}
\setlength{\parskip}{4pt}
\noindent\textbf{Motivation.} Weighted gene co-expression network analysis (WGCNA) and its
variants remain the dominant paradigm for partitioning genes into functional modules.
We treat the validity of this paradigm as a statistical question rather than a modelling
choice, and test it on a fixed benchmark with strictly paired comparisons and null models
that separate genuine biological organisation from artefacts of network construction. Four
assumptions are examined: that a monotone co-expression metric captures functional
relatedness; that a global cut yields reproducible modules; that a single partition resolves
nested sub-structure; and that reported modularity reflects biological organisation. All
four fail.

\vspace{4pt}
\noindent\textbf{Results.} We assembled a benchmark of 16 publicly available expression
cohorts spanning neoplastic and non-neoplastic disease, nine GO biological-process
pathways and 31 partitioning algorithms from eight methodological families, yielding 100
pathway--cohort pairs and 3100 evaluation records. All algorithms operated on an identical
topological overlap matrix per pair, enforcing a strictly paired comparison. Evaluation
followed a six-layer protocol covering structural compactness, functional coherence
(GO/KEGG/Reactome-referenced, size-unbiased $z$-scores), bootstrap stability,
multi-perturbation robustness, biological validity (curated antagonistic gene pairs) and
downstream clinical prediction. To eliminate circularity, graph representation learning was
reformulated as self-supervised edge prediction, removing all community labels from
training. Synthetic benchmarks with known ground truth supplied absolute accuracy
references. Motivated by the finding that expression graphs degrade under realistic
constraints, we further introduce \textbf{FACEO-R}, a functionally anchored overlapping
community detection method in which the weight of GO-derived functional similarity adapts
to the estimated reliability of the expression graph. Under controlled degradation (sample
subsampling, expression noise and feature dropout) and on clean data alike, FACEO-R
outperformed Leiden at $\gamma = 1.0$ in 118/145 comparisons (81\%;
paired Wilcoxon $p = 1.9\times10^{-14}$) and WGCNA-style
hierarchical cutting in 127/145 (88\%; $p = 9.0\times10^{-20}$), with
the largest gains under feature dropout ($\Delta = +0.070$, 27/28 wins). Multi-resolution community detection produced functionally more coherent
modules than hierarchical clustering (median $z = 0.689$ versus $0.045$; paired Wilcoxon
$q = 4.6\times10^{-7}$) while simultaneously improving modularity
($0.291$ versus $0.000$; $q = 2.0\times10^{-17}$), and recovered nested functional
sub-structures with markedly higher accuracy in synthetic benchmarks (ARI $1.000$ versus
$0.908$). Downstream clinical prediction also favoured multi-resolution partitions
(AUC $0.989$ versus $0.898$).

\vspace{4pt}
\noindent\textbf{Conclusions.} We find that the three implicit assumptions of current
practice --- a single adequate algorithm, a single defensible resolution, and modularity as
evidence --- all fail, and that the failures are systematic. First,
modularity --- and its significance against degree-preserving null models --- proved unstable
with respect to graph construction (threshold sensitivity reverses the sign; weights
contribute 29\%), so it should not be used as primary evidence. Second, the functionally optimal
resolution is \emph{not predictable}: optimal $\gamma$ was highly dispersed across pathways
(entropy 1.74 of a maximum 1.79, with no value optimal in more than 26\% of cases), yet four
families of automatic selection strategies --- structural-fingerprint regression, pathway-size
grouping, multi-resolution consensus clustering and cross-library held-out annotation
guidance --- all failed (accuracy 0.16--0.46, none significantly better than chance, with
consensus clustering actively harmful). Third, any fixed resolution therefore sacrifices a
mean of 0.765 functional-coherence $z$-units relative to per-pathway optimal selection. We
recommend reporting resolution-sensitivity curves validated against held-out functional
evidence, rather than modularity, module counts, or a single arbitrary resolution.

\vspace{6pt}
\noindent\textbf{Availability and implementation:} The FACEO-R implementation, all
analysis scripts, processed expression matrices and complete evaluation tables are freely
available at \url{https://github.com/sunmy1129/faceo-r} under an MIT licence.\\
\noindent\textbf{Contact:} weixiong.zhang@polyu.edu.hk; mingyu.sun@connect.polyu.hk\\
\noindent\textbf{Keywords:} functional module partitioning; community detection; Leiden
algorithm; WGCNA; graph neural networks; self-supervised learning; benchmarking
\end{minipage}
\vspace{0.30em}
\hrule height 0.4pt
\vspace{0.85em}
}]

\vspace{2pt}

%% ============================================================

\renewcommand{\thefigure}{S\arabic{figure}}
\setcounter{figure}{0}
\section*{Appendix}
%% ============================================================

\begin{figure}[H]
\centering
\includegraphics[width=\columnwidth,height=0.30\textheight,keepaspectratio]{../output/figures_v3/FigCA_clean.png}
\caption{The advantage is present with and without degradation, and the anchoring strength
has a narrow optimum. (a) Mean functional coherence under clean data and five degradation
regimes; FACEO-R exceeds both baselines in every condition. (b) Coherence as a function of a
fixed anchoring weight $\alpha$ (solid); performance improves monotonically as anchoring
strengthens, the originally proposed adaptive schedule under-anchored (dashed, below every
fixed setting in $[0.05, 0.35]$), and the retuned schedule matches the best fixed weight.}
\label{fig:clean}
\end{figure}

%% ============================================================





\begin{figure}[H]
\centering
\includegraphics[width=\columnwidth,height=0.28\textheight,keepaspectratio]{../output/figures_v3/FigCA_bio.png}
\caption{Recovered modules map onto coherent, specific biological processes. Median
enrichment significance of the top Gene Ontology term per module, with the input pathway
excluded from consideration; annotations show the median number of module genes in each
term. Twenty-two distinct processes were recovered across four cohorts.}
\label{fig:bio}
\end{figure}

\end{document}

\end{document}
